\documentclass[11pt]{article}
\usepackage[T1]{fontenc}
\usepackage[letterpaper,margin=0.88in]{geometry}
\usepackage{amsmath,amssymb,amsthm}
\usepackage{graphicx,microtype}
\usepackage[hidelinks]{hyperref}

\newtheorem*{theorem}{Theorem}
\newtheorem*{corollary}{Corollary for Open Problem (I)}
\newcommand{\E}{\mathbb E}
\newcommand{\piu}{\pi}

\title{The zero endpoint is isolated from every positive normalized
$r$-logSob index}
\author{A partial result for Open Problem (I) in arXiv:1108.1210}
\date{21 August 2026}

\begin{document}
\maketitle

\section*{Source problem}
Mossel, Oleszkiewicz, and Sen ask for which intervals $I$ the normalized
$r$-logSob inequalities are uniformly equivalent over all finite reversible
Markov semigroups: a constant at one index should imply a constant at every
other index with loss depending only on $I$ \cite[Open Problem (I),
Section~12]{MOS}.

\begin{center}
\includegraphics[width=0.96\textwidth]{figures/open_problem_crop.png}

{\small Open Problem (I), official arXiv PDF page 27.}
\end{center}

The source itself separates indices $0$ and $1$ using bounded-degree
expanders.  The result below fills the whole punctured interval between them
on the zero-endpoint side.

\section*{Proof Intuition}
An $r$-logSob inequality with fixed $r>0$ gives reverse
hypercontractivity after a time proportional to its constant.  Testing the
resulting two-function inequality on two single vertices forces a polynomial
lower bound on every transition probability.  In a bounded-degree graph,
some vertices are logarithmically far apart, while a continuous-time walk can
reach them before that many Poisson-distributed jumps only with a much smaller
probability.  Therefore the $r$-logSob constant must grow at least like
$\log n$.  Expanders keep the $0$-logSob (Poincare) constant bounded, so zero
cannot share an equivalence interval with any positive index.

\begin{theorem}
Fix $0<r<1$ and an integer $d\ge3$.  There are constants
$c_{r,d}>0$ and $N_{r,d}$ such that the following holds.  If $G$ is a
connected $d$-regular graph on $n\ge N_{r,d}$ vertices, $P$ is its
simple-random-walk kernel, and the semigroup generated by $L=I-P$ satisfies
the normalized $r$-logSob inequality with constant $C$, then
\[
                         C\ge c_{r,d}\log n.
\]
\end{theorem}

\begin{proof}
Put $s=r/2\in(0,1)$ and $s'=s/(s-1)<0$.  Proposition~6.1 of
\cite{MOS}, applied with output exponent $s'$ and input exponent $s$, gives
\[
       \|T_t h\|_{s'}\ge \|h\|_s,
       \qquad t=\kappa_r C,
       \qquad \kappa_r=-\tfrac12\log(1-s)>0.
\]
Indeed the proposition's time is
\[
 \frac C4\log\frac{1-s'}{1-s}
 =-\frac C2\log(1-s).
\]
Reverse Holder, followed by the displayed reverse-hypercontractive estimate,
therefore yields for nonnegative $f,g$
\[
             \E_{\piu}[fT_tg]
       \ge \|f\|_s\,\|T_tg\|_{s'}
       \ge \|f\|_s\,\|g\|_s.
\]
Take $f=\mathbf1_{\{u\}}$ and $g=\mathbf1_{\{v\}}$.  Since the stationary
measure is uniform,
\[
 \frac1n P_t(u,v)\ge n^{-2/s},
 \qquad\text{hence}\qquad
 P_t(u,v)\ge n^{-A_r},\quad A_r=\frac4r-1.
 \tag{1}
\]

A ball of radius $R$ in a graph of maximum degree $d$ has at most
$1+d((d-1)^R-1)/(d-2)$ vertices.  Consequently, for all sufficiently large
$n$ there are $u,v$ at distance
\[
                         D\ge a_d\log n
\]
with a constant $a_d>0$ depending only on $d$.  The continuous-time walk
makes a Poisson$(t)$ number $J$ of jumps.  It cannot reach $v$ from $u$ in
fewer than $D$ jumps, and for $t<D$ the Chernoff bound gives
\[
             P_t(u,v)\le \Pr\{J\ge D\}
                    \le \left(\frac{et}{D}\right)^D.       \tag{2}
\]
Choose $0<\theta_{r,d}<e^{-1}$ so small that
$a_d\log(e\theta_{r,d})<-2A_r$.  If $t\le\theta_{r,d}D$,
then (2) is at most $n^{-2A_r}$, contradicting (1).  Thus
\[
 \kappa_r C=t>\theta_{r,d}D
   \ge \theta_{r,d}a_d\log n,
\]
which proves the theorem with
$c_{r,d}=\theta_{r,d}a_d/\kappa_r$.
\end{proof}

\begin{corollary}
No interval $I\subset[0,2]$ containing both $0$ and a positive number has
the equivalence property in Open Problem~(I).
\end{corollary}

\begin{proof}
For a fixed bounded-degree spectral expander family, the Poincare constants
are uniformly bounded.  Lemma~3.1 of \cite{MOS} identifies this, up to its
stated factor, with a uniform $0$-logSob constant.  For any fixed
$r\in I\cap(0,1)$ the theorem makes the optimal $r$-logSob constants diverge.
If the positive point of $I$ is at least one, interval convexity supplies a
point in $(0,1)$; alternatively the source already supplies the $r=1$
separation.  Hence uniform equivalence is impossible.
\end{proof}

\section*{Verification, limitations, and novelty}
The exponent conversion, indicator normalization, bounded-degree diameter
bound, and Poisson Chernoff comparison were independently rederived from the
official source PDF.  The algebra regression in
\texttt{code/verify\_expander\_algebra.py} checks the chosen constants on a
range of indices and degrees; it is not used as proof.

This settles only intervals that contain zero.  It does not compare two fixed
indices in $(0,1)$ and does not classify the intervals bounded away from zero.

Cheap run indexes plus bounded exact-phrase, ``$r$-logSob expander'',
positive-index, and reverse-hypercontractivity searches through 21 August
2026 found the source and later work on other reverse-hypercontractive
settings, but no claimed resolution of this zero-versus-every-positive-index
statement.  Novelty remains subject to human review.

\paragraph{Human-review focus.}
Check the admissible exponent range in source Proposition~6.1, the factor
$\kappa_r$, and the conversion from stationary expectation to the transition
probability in (1).  The remainder is the standard bounded-degree ball count
and Poisson Chernoff bound.

\begin{thebibliography}{9}
\bibitem{MOS}
E. Mossel, K. Oleszkiewicz, and A. Sen,
\emph{On reverse hypercontractivity}, arXiv:1108.1210v2 (2018),
especially Proposition~6.1 and Section~12.
\end{thebibliography}

\end{document}
